\section{Results}
\label{sec:results}

\subsection{HistoAtlas turns routine pathology into a named, compartment-resolved measurement layer}
\label{sec:results:resource}

We built HistoAtlas from \num{6745} H\&E-stained diagnostic whole-slide images
from \num{6088} patients spanning 21 TCGA solid-tumor types
(Table~S1). Twelve further TCGA types were excluded because their dominant cell
morphologies fall outside the training domain of the cell detection model,
including lymphoid, glial, melanocytic, mesenchymal, neuroendocrine, renal
tubular and germ-cell tumors (Table~S2). Descriptive analyses of the measurement
space use every slide; all statistical associations use one representative slide
per patient, so that the 312 patients with several slides are not counted more
than once.

The pipeline converts each slide into quantitative morphology in two stages
(Figure~\ref{fig:fig1}A). A tissue model assigns every region to one of five
compartments (cancerous epithelium, stroma, necrosis, normal epithelium, blood);
tumor and stroma together account for more than 90\% of analyzed tissue area, so
downstream measurements can be expressed relative to biologically meaningful
territory rather than to whole-slide area. The HistoPLUS cell
model~\cite{adjadj2026histoplus}, which classifies cells into 13 types, then
detects individual cells, of which nine types are used here (tumor cell,
lymphocyte, fibroblast, plasmocyte, neutrophil, eosinophil, red blood cell,
apoptotic body, mitotic figure), more
than four billion cells across the cohort. Signed distance transforms from the
tumor boundary define spatial bands, so that a lymphocyte counted inside tumor
epithelium is not confused with one counted 150~\textmu m away in stroma. From
these maps we derived 38 histomic measurements in six groups: tissue composition
(3), cell densities (8), nuclear morphology and kinetics (7), spatial organization
(11), spatial heterogeneity (5) and cell-type ratios (4) (Table~S3).

Across 22 analysis cohorts (21 cancer types
plus a pan-cancer set) we fitted \num{5623} Cox proportional hazards models
(four endpoints, each under an unadjusted and a covariate-adjusted
specification), of which \num{5283} yielded valid estimates after
proportional-hazards screening, and computed \num{487638} Spearman correlations
against messenger RNA expression, pathway activity, copy number and immune
deconvolution scores. Benjamini--Hochberg correction was applied within declared,
non-overlapping families (150 for correlations, defined by cancer type, molecular
data type and model). At a false discovery rate of 0.05, 260 survival
associations and \num{88920} molecular correlations were significant
(\num{87217} of them involving the 38 measurements; the others involve three
quality-control variables released in the same table; Table~S8). Every association carries an effect size, a bootstrap 95\% confidence
interval, raw and adjusted $P$ values, a minimum detectable effect at 80\% power,
a proportional-hazards flag where applicable, and the identifier and size of its
correction family.

Pairwise Spearman
correlation across the 38 measurements resolves coherent modules: immune-density
measurements form a tight positive block, nuclear morphometry clusters
separately, and cross-module anticorrelations, such as between tumor area
fraction and the tumor--stroma interface measurements, delineate distinct axes
(Figure~\ref{fig:fig1}B). Principal component analysis gives an effective
dimensionality of 18 components at 90\% explained variance, so the layer is
compressible but not collinear. A UMAP embedding of all \num{6745} slides
recovers expected morphological neighborhoods, with breast, liver, thyroid and
endometrial carcinoma occupying distinct regions and squamous and adenocarcinoma
types interleaved in the center (Figure~\ref{fig:fig1}C). Unsupervised $K$-means
partitioning into ten pan-cancer clusters recovers groups whose cancer-type
composition and measurement profiles are interpretable without any molecular
input (Figure~\ref{fig:fig1}D,E), and which are enriched for distinct immune
subtypes and Hallmark pathway states when molecular annotation is subsequently
overlaid (Figure~S1). Two of these purely morphological groups separate
prognostically in Cox models stratified by cancer type, so the effects are not
restatements of tissue of origin: a low-proliferation cluster dominated by liver,
thyroid and adrenocortical carcinoma (C2) was associated with better overall
survival (HR $= 0.54$, 95\%~CI $[0.40, 0.73]$, $P_{\mathrm{adj}} = 6.3 \times
10^{-4}$; 516 patients in the cluster, 108 events), and a high-interface cluster
dominated by breast and prostate carcinoma (C8) with worse survival (HR $= 1.36$,
95\%~CI $[1.15, 1.62]$, $P_{\mathrm{adj}} = 1.8 \times 10^{-3}$; \num{1112}
patients, 214 events).

\subsection{Six continuous morphological programs summarize the measurement space}
\label{sec:results:programs}

We applied
sparse principal component analysis ($\alpha = 1.0$) to the 38 measurements
across \num{5371} TCGA patients with complete measurement vectors, yielding six
continuous programs, each loading on 29 to 33 measurements with a few dominant
ones (Figure~\ref{fig:fig2}A). Their biological identities were not assigned in
advance but read off from correlations with 50 MSigDB Hallmark pathway scores~\cite{liberzon2015}
(Figure~\ref{fig:fig2}C): program~1 correlates most strongly with
\textsc{allograft\_rejection} ($\rho = 0.59$, $n = \num{5371}$) and the
interferon-$\gamma$ response ($\rho = 0.51$), program~2 with
\textsc{estrogen\_response\_early} ($\rho = 0.44$) and
\textsc{epithelial\_mesenchymal\_transition} ($\rho = 0.41$), and program~3 with
\textsc{myc\_targets\_v1} and \textsc{g2m\_checkpoint} (both $\rho = 0.53$). We
therefore label them immune infiltration, stromal architecture, nuclear atypia
and proliferation, tumor density, myeloid inflammation and fibroblast biology.
The same correspondence is
visible at the level of individual measurements, where immune densities align
with immune pathway signatures, nuclear morphometry and mitotic index with cell
cycle and proliferation signatures, and stromal and interface measurements with
the epithelial--mesenchymal transition, angiogenesis and coagulation signatures
(Figure~S2A).

\subsection{Compartment-resolved immune topology is prognostic and molecularly anchored}
\label{sec:results:immune}

Tumor-infiltrating lymphocytes are conventionally scored in the stromal
compartment~\cite{salgado2015til}; HistoAtlas measures lymphocyte density
separately in tumor epithelium and in stroma. Intratumoral
lymphocyte density was associated with improved overall survival in the
pan-cancer adjusted model (hazard ratio [HR] $= 0.87$, 95\%~CI $[0.81, 0.93]$,
$P_{\mathrm{adj}} = 9.8 \times 10^{-4}$, $n = \num{4560}$, \num{1356} events),
with the strongest cancer-specific effects in breast carcinoma (HR $= 0.72$,
95\%~CI $[0.60, 0.88]$, $P_{\mathrm{adj}} = 0.018$, $n = 960$) and head-and-neck
squamous carcinoma (HR $= 0.74$, 95\%~CI $[0.63, 0.87]$,
$P_{\mathrm{adj}} = 3.9 \times 10^{-3}$, $n = 444$; Figure~\ref{fig:fig3}A,B).
Stromal lymphocyte density, measured on the same slides by the same model,
showed a weaker pan-cancer effect (HR $= 0.89$, 95\%~CI $[0.83, 0.97]$,
$P_{\mathrm{adj}} = 0.031$, $n = \num{4561}$) and no significant effect in breast
carcinoma (HR $= 0.93$, 95\%~CI $[0.77, 1.12]$, $P_{\mathrm{adj}} = 0.67$). The
two compartments therefore separate: localization within tumor epithelium, not
lymphocyte abundance alone, carries the stronger prognostic
signal, consistent with the prognostic value of immune-cell location reported in
colorectal cancer~\cite{galon2006, pages2018}. In most of the 18 cancer types with an
evaluable adjusted model the intervals are wide and include 1, which is the expected
behavior of small cohorts (Figure~S3B).

The same distinction appears in the molecular data. Tumor--lymphocyte
nearest-neighbor distance at the invasive front, a direct spatial measure of
immune exclusion~\cite{joyce2015exclusion}, was inversely associated with
\textit{CD8A} expression in breast carcinoma ($\rho = -0.53$, 95\%~CI
$[-0.58, -0.48]$, $P_{\mathrm{adj}} = 1.8 \times 10^{-68}$, $n = 958$;
Figure~\ref{fig:fig3}C), and the top gene correlates of intratumoral lymphocyte
density are dominated by T-cell, B-cell and checkpoint transcripts
(\textit{TIGIT}, \textit{CD3E}, \textit{CD3D}, \textit{CD79A}, \textit{PDCD1},
\textit{CD8A}, \textit{CTLA4}), with luminal markers (\textit{ESR1},
\textit{PGR}, \textit{CCND1}) among the strongest negative correlates
(Figure~\ref{fig:fig3}D).

\subsection{Named measurements recover proliferation, invasion and differentiation programs}
\label{sec:results:bridge}

In TCGA breast
carcinoma, mitotic index correlated with canonical cell-cycle regulators,
including \textit{CCNE1} ($\rho = 0.568$, 95\%~CI $[0.519, 0.617]$,
$P_{\mathrm{adj}} = 2.7 \times 10^{-80}$, $n = 958$), \textit{PLK1}
($\rho = 0.558$ $[0.510, 0.607]$, $P_{\mathrm{adj}} = 5.1 \times 10^{-77}$),
\textit{AURKA} ($\rho = 0.539$, $P_{\mathrm{adj}} = 4.7 \times 10^{-71}$),
\textit{CHEK1} ($\rho = 0.526$) and \textit{BIRC5} ($\rho = 0.520$; all
$P_{\mathrm{adj}} < 10^{-64}$). Invasion depth, defined as the 75th percentile of
tumor distance from the tissue boundary, aligned with an
epithelial--mesenchymal axis (\textit{ZEB1}: $\rho = 0.318$
$[0.262, 0.375]$, $P_{\mathrm{adj}} = 8.4 \times 10^{-23}$, $n = 957$;
\textit{ALDH1A1}: $\rho = 0.327$, $P_{\mathrm{adj}} = 3.8 \times 10^{-24}$).
Nuclear pleomorphism, the coefficient of variation of tumor nuclear area,
correlated positively with proliferation genes (\textit{PLK1}:
$\rho = 0.484$, $P_{\mathrm{adj}} = 1.4 \times 10^{-55}$) and negatively with
luminal differentiation markers (\textit{BCL2}: $\rho = -0.374$,
$P_{\mathrm{adj}} = 8.1 \times 10^{-32}$; \textit{ESR1}: $\rho = -0.358$,
$P_{\mathrm{adj}} = 4.4 \times 10^{-29}$). These correspondences restate, in
continuous and automatically measured form, relationships that underpin
histopathological grading~\cite{elston1991grading}. Their value here
is that they are computed identically on every slide in the atlas, so they can
be compared across 21 cancer types rather than asserted one tumor type at a time.
Across all molecular data types the coupling is strongest for expression and
pathway activity and weakest for copy number: 80\% of pan-cancer adjusted
expression correlations and 83\% of pathway correlations were significant,
against 52\% for copy number (Figure~S2B).

\subsection{The measurement layer and its molecular links transfer to an independent consortium}
\label{sec:results:transfer}

Portability was assessed on \num{1095} CPTAC slides from 817 patients across five
cancer types with TCGA counterparts: breast ($n = 198$ slides), lung
adenocarcinoma ($n = 240$), head-and-neck squamous carcinoma (390 slides from 112
patients), pancreatic adenocarcinoma ($n = 175$) and endometrial carcinoma
($n = 92$).

For each cancer type we computed the standardized difference between cohort
medians, in units of the TCGA standard deviation, for each measurement, with a
prespecified criterion of at least 31 of 38 measurements within one standard
deviation. Four of five types passed: lung adenocarcinoma and pancreatic
adenocarcinoma at 34/38 (89.5\%), head-and-neck at 32/38 (84.2\%) and breast at
31/38 (81.6\%); endometrial carcinoma narrowly missed at 30/38 (78.9\%), driven by
tissue-composition shifts in the smallest cohort ($n = 92$ versus $n = 459$ in
TCGA; Figure~\ref{fig:fig4}A; Table~S9). The inter-measurement correlation
geometry was preserved across cohorts (Pearson $r = 0.84$ over the 703
upper-triangle pairs, $P < 10^{-180}$; Table~S9), indicating that the 38
measurements span the same latent morphological space in independent tissue.

The morphology-to-transcriptome bridge replicated. All ten prespecified
histomic--molecular associations matched in direction and nine remained
significant after correction within the full CPTAC family of approximately
\num{24700} tests per model per cancer type, including mitotic index with
\textit{PLK1} in lung adenocarcinoma ($\rho = 0.505$,
$P_{\mathrm{adj}} = 5.0 \times 10^{-5}$), invasion depth with \textit{ZEB1} in
breast carcinoma ($\rho = 0.319$, $P_{\mathrm{adj}} = 4.6 \times 10^{-3}$) and
nuclear pleomorphism with \textit{ESR1} ($\rho = -0.389$,
$P_{\mathrm{adj}} = 2.7 \times 10^{-4}$) and \textit{BCL2}
($\rho = -0.416$, $P_{\mathrm{adj}} = 7.7 \times 10^{-5}$). Beyond the
prespecified set, effect sizes were concordant at scale. Across 750 matched
feature--gene--cancer triples drawn from each cohort's strongest correlates, the
two cohorts' Spearman coefficients correlated at Pearson $r = 0.787$ (95\%~CI
$[0.758, 0.813]$, $P = 4.0 \times 10^{-159}$), with mean $|\rho|$ modestly
attenuated in CPTAC (0.251 versus 0.301). Across all \num{25155} evaluable
measurement--gene pairs, most of which are near zero in both cohorts, the
correlation was $r = 0.56$ (Figure~\ref{fig:fig4}B), as expected when weak
associations are estimated in small cohorts.

Program transfer was mixed in an informative way. Projecting the six
TCGA-fitted programs into 297 CPTAC patients with complete measurement vectors,
two-sample Kolmogorov--Smirnov tests showed no significant distributional shift
for the immune program ($D = 0.072$, $P = 0.10$) or the fibroblast program
($D = 0.064$, $P = 0.19$), a borderline shift for the stromal program
($D = 0.096$, $P = 0.010$), and larger shifts for the remaining three
($D = 0.15$ to $0.34$, $P < 10^{-4}$; Figure~\ref{fig:fig2}B). The programs also
validated at the protein level: the immune program tracked CXCL9
($\rho = 0.85$), CD79A ($\rho = 0.82$) and CD3E ($\rho = 0.78$) protein
abundance, and the stromal program tracked ZEB1 ($\rho = 0.72$) and FN1
($\rho = 0.71$), across \num{2220} program--protein pairs in four cancer types
(Table~S10). By contrast, transferring the ten hard cluster labels by
nearest-centroid assignment reproduced pathway enrichment only weakly (Spearman
$\rho = 0.15$, $P = 1.4 \times 10^{-3}$ across 450 cluster--pathway pairs), with
three of nine evaluable clusters concordant ($\rho = 0.38$ to $0.51$) and two
inverted ($\rho = -0.59$ and $-0.65$; Figure~S4). Adjusted Rand index between
cluster labels and cancer type was low in both cohorts (0.14 in TCGA, 0.07 in
CPTAC), confirming that the clusters capture morphology beyond tissue of origin.

\subsection{The bridge extends to protein, and the information added beyond bulk omics is incremental}
\label{sec:results:protein}

CPTAC's matched proteomics allowed the bridge to be tested one layer further
out. Across \num{14687} measurement--gene pairs with both messenger RNA and protein
measured, the two modalities' coefficients correlated at $r = 0.63$
(Figure~\ref{fig:fig4}C). Among the 320 pairs significant in both, directional
concordance was 98.8\% (316/320, 95\%~CI $[96.8\%, 99.7\%]$), and the median
protein-minus-RNA difference in effect size was essentially zero
($\Delta\rho = -0.007$). Protein was therefore not globally stronger than RNA,
but it was specifically stronger for markers where post-translational biology
dominates: nuclear pleomorphism tracked ESR1 protein more closely than
\textit{ESR1} transcript in breast carcinoma ($\rho = -0.440$ versus $-0.389$),
consistent with estrogen receptor protein being the phenotypic determinant of
luminal differentiation~\cite{allison2020erpr}, and stromal
measurements tracked vimentin protein more closely than transcript
($\rho = 0.578$ versus $0.344$). Mitotic index tracked cell-cycle proteins in
lung adenocarcinoma (MKI67, CDK1, TOP2A, CHEK1, PLK1), and intratumoral lymphocyte
density tracked CD79A, CD3E, GZMA, PRF1, GZMB and IDO1 protein in breast
carcinoma. Of the four biological axes probed by this triangulation, proliferation markers
are supported in all four lines of evidence, immune and hormonal markers in three
or four, and epithelial--mesenchymal markers mainly in TCGA transcripts and CPTAC
protein (Figure~\ref{fig:fig4}D).

Because protein abundance is largely set by transcript abundance, we asked how
much of this signal is protein-specific. Partial Spearman correlations between
each measurement and each protein, adjusting for the matched transcript, left
132 of \num{14687} triplets significant after correction (0.9\%), 120 of them in
breast carcinoma, including stromal and DNA-repair markers such as VIM
(Figure~\ref{fig:fig5}C); 77\% of protein correlations retained more than half
their magnitude after adjustment. Immune proteins such as CD8A and GZMB were
largely absorbed.

The same question can be asked of the survival results, and gives a
correspondingly measured answer. For each of three measurements we refitted the
strongest clinical Cox model with matched molecular covariates added
(\textit{CD8A}, \textit{TIGIT}, immune subtype and leukocyte fraction for the
immune measurements; the \textsc{g2m\_checkpoint} score and \textit{MKI67} for
mitotic index; Figure~\ref{fig:fig5}A). The hazard ratio for intratumoral
lymphocyte density in head-and-neck carcinoma was essentially unchanged (0.76 to
0.72, still significant), whereas the hazard ratios for tumor--lymphocyte
distance in cervical carcinoma (1.34 to 1.18) and mitotic index in mesothelioma
(2.00 to 1.40) moved toward 1 and were no longer significant, so part of those
associations is carried by the matched transcripts. Adding a measurement to a clinical-plus-RNA model raised Harrell's
apparent (in-sample) concordance index by a median of $+0.0007$ to $+0.0027$ across the 18 evaluable
cancer types for the three pan-cancer measurements ($+0.0038$ for nuclear
pleomorphism, tested in breast carcinoma only), and the increment was positive in
only 50\% to 61\% of individual cancer types; in
breast carcinoma, the largest cohort, intratumoral lymphocyte density added
$+0.008$ over clinical covariates plus \textit{CD8A} (Figure~\ref{fig:fig5}B).

\subsection{Calibration and sensitivity analyses mark where the atlas is reliable}
\label{sec:results:calibration}

Because HistoAtlas is meant to be reused, we characterized where its outputs are
stable and where they are not. Under a permutation null in which molecular
labels were shuffled within cancer type and the full production correction
pipeline was re-run (100 permutations over the \num{39900} measurement $\times$
pathway $\times$ cancer-type tests), the mean null discovery rate was
$1 \times 10^{-6}$ against an observed rate of 26.7\%
(empirical $P = 0.0099$, the smallest attainable with 100 permutations), so the
discovery rate is not a trivial consequence of sample size. Replacing
Benjamini--Hochberg with the more conservative Benjamini--Yekutieli procedure,
which permits arbitrary dependence among tests, retained 68.1\% of significant
correlations and 54.0\% of significant survival associations under a single
pooled correction family (Table~S11).

Targeted sensitivity analyses locate the confounders that matter. Adjusting for
molecular subtype attenuated the leading histomic--molecular correlations by
4.3\% in median $|\rho|$ in breast carcinoma and 6.2\% in endometrial carcinoma,
but by 28\% in pancreatic adenocarcinoma, where immune subtype carries much more
of the signal (Table~S12). Adjusting survival models for leukocyte fraction did
not erase the leading prognostic associations; slightly more measurements reached
significance afterwards (28/179 versus 25/179), indicating that they are not
simply surrogates for immune content (Table~S13). Batch effects from tissue
source site were present but not dominant: pan-cancer variance decomposition
attributed 44.7\% of variance to source site, 32.7\% to cancer type and 22.6\% to
residual, while within the 20 cancer types with several sites, source site
accounted for a median of 21.3\%, and all per-cancer silhouette scores by source
site were negative (Figure~S3A). Statistical power is the binding constraint in
small cohorts: breast carcinoma can detect HR~$\geq 1.62$ at 80\% power, whereas
cholangiocarcinoma ($n = 36$) requires HR~$\geq 3.75$, which is why \num{4846} of
\num{5623} survival associations carry an \emph{insufficient} evidence badge
(Figure~S3B).

Some measurements are more protocol-sensitive than others. Median tumor nuclear
irregularity shifted systematically between cohorts in all five matched cancer
types (standardized difference 1.8 to 2.4 TCGA standard deviations, largest in
lung adenocarcinoma; Figure~\ref{fig:fig4}A), a scale of shift that would
distort any cross-cohort comparison of that measurement. The atlas flags it
accordingly. Taken together, these analyses argue that HistoAtlas should be
consumed as a calibrated resource rather than an undifferentiated catalog: the
measurement layer is portable, the continuous programs are more reliable
cross-cohort units than the hard clusters, and the strongest claims are those
supported simultaneously by TCGA discovery, CPTAC replication and sensitivity
analyses that bound the plausible confounders.

\subsection{An open atlas makes pan-cancer morphology reusable by people and by machines}
\label{sec:results:atlas}

The statistical layer described above is released, not summarized. HistoAtlas is
served as a static web application at
\href{https://histoatlas.com}{histoatlas.com}, backed by precomputed JSON so
that no analysis server is required at runtime, and as a machine-queryable
interface exposing 45 typed tools over the same tables. A user can enter from
any of four directions: a measurement (what does intratumoral lymphocyte density
predict, and where), a cancer type (which measurements matter here), a molecular
entity (which morphologies track \textit{ESR1}), or a slide (what does this case
look like relative to the cohort). Each association is displayed with its effect
size, confidence interval, adjusted $P$ value, evidence badge, minimum
detectable effect, correction-family size and proportional-hazards status, so
that a weak result is legible as weak rather than merely present. Measurement
tables, per-slide values, cluster assignments and the full association tables
are downloadable.

Selecting a survival association leads to the measurement's definition
(Figure~\ref{fig:fig6}), from
there to the per-slide compartment map on which it was computed, and from there
to the individual image tiles that contributed most to the value, overlaid with
the cell-level predictions that generated them. The contrast between
intratumoral and stromal lymphocyte density in breast carcinoma is, at the top
of the stack, two hazard ratios (Figure~\ref{fig:fig6}B); at the bottom it is
the tiles that drive each score on one slide (Figure~\ref{fig:fig6}C,D). The
same view also exposes what a number hides: the third-ranked intratumoral tile on
this slide is dominated by fibroblasts rather than lymphocytes, a case a user can
see, question and report. A user who distrusts a number can look at the cells
behind it, which is not possible for a resource built on learned embeddings.
The programmatic interface exposes the same objects as structured tool calls,
so the atlas can also be queried by automated agents conducting systematic
hypothesis generation rather than only by humans browsing pages.
